\begin{table}[ht]
\centering
\footnotesize
\renewcommand{\arraystretch}{1.3}
\caption{Confound taxonomy and mitigation status for the pilot study.}
\label{tbl:confound-taxonomy}
\begin{tabular}{@{}p{3.2cm}p{5.5cm}p{3.8cm}p{2cm}@{}}
\toprule
\textbf{Confound Type} & \textbf{Description} & \textbf{Pilot Status} & \textbf{Mitigation Stage} \\
\midrule
\textbf{Run-order} & All Condition~B runs preceded Condition~A. Time-of-day, experimenter familiarity, or tool snapshots may contribute to observed contrast. & Fully confounded & Validation \\
\addlinespace[0.5em]
\textbf{Prompt verbosity} & Condition~A prompt is longer (631 tokens) than Condition~B (196 tokens). Increased verbosity may improve output quality independent of convention encoding. & Uncontrolled & Validation \\
\addlinespace[0.5em]
\textbf{Service integration} & SmartKit integrates billing, auth, and email. If the pilot feature called these services, pattern-matching from integration code could confound convention adherence measurement. & Eliminated by design & --- \\
\addlinespace[0.5em]
\textbf{Component-level mechanism} & AGENTS.md, FSD layout, and sibling Zod patterns are bundled in Condition~A. The pilot measures ecosystem-level effect but cannot assign contribution to individual components. & Bundled (intentional) & Future Work (item~4) \\
\addlinespace[0.5em]
\textbf{Sample size} & $n=5$ per condition is sufficient for deterministic-verdict detection but insufficient for correlation analysis or distribution comparison. & Design constraint & Validation \\
\bottomrule
\end{tabular}
\end{table}
